% Part: first-order-logic
% Chapter: introduction
% Section: models-theories

\documentclass[../../../include/open-logic-section]{subfiles}

\begin{document}

\olfileid[hi]{fol}{int}{mod}

\olsection{प्रतिरूप और सिद्धान्त}

प्रथम-क्रम तर्क का वाक्यविन्यास और अर्थविज्ञान परिभाषित कर लेने पर
!!{structure}s हमारे सामने होती हैं और हम उनके तथा अर्थवैज्ञानिक धारणाओं के गुणों की जाँच आरंभ कर सकते हैं।
हम !!{derivation} प्रणालियाँ भी परिभाषित करके उनकी जाँच कर सकते हैं।
!!{sentence}s के किसी समुच्चय के बारे में हम पूछ सकते हैं: कौन-सी
!!{structure}s उस समुच्चय के सभी !!{sentence}s को सत्य बनाती हैं?
!!{sentence}s का समुच्चय~$\Gamma$ दिया हो, तो उन्हें संतुष्ट करने वाली
!!a{structure}~$\Struct{M}$ को \emph{~$\Gamma$ का प्रतिरूप} कहते हैं। हम
~$\Gamma$ से आरंभ करके उसके प्रतिरूप खोजने का प्रयत्न कर सकते हैं---वे कैसे
दिखते हैं? उन्हें कितना बड़ा या छोटा होना पड़ता है? किंतु हम किसी एक
!!{structure} से अथवा ऐसे किसी संग्रह से भी आरंभ कर सकते हैं जिसमें
!!{structure}s हों, और यह पूछ सकते हैं: उसमें किन !!{sentence}s को सत्य माना
जाता है? क्या उन !!{sentence}s में कुछ ऐसे हैं जिनके द्वारा ठीक ये
!!{structure}s इस अर्थ में \emph{लक्षणित} होती हों कि वे वाक्य इन्हीं में,
और केवल इन्हीं में, सत्य हों? ऐसे प्रश्न \emph{प्रतिरूप सिद्धान्त} के विषय
हैं। वे \emph{अभिगृहीतीय विधि} के आधार में भी हैं: ऐसे किसी संग्रह का वर्णन
करना जिसमें !!{structure}s हों, और उसका वर्णन !!{sentence}s के ऐसे समुच्चय
से करना जो किसी सिद्धान्त के अभिगृहीत हों। यह इस अवलोकन से संभव होता है कि
अभिगृहीतों से प्रथम-क्रम तर्क में जिन !!{sentence}s की निष्पत्ति होती है,
ठीक वे सभी अभिगृहीतों के प्रतिरूपों में सत्य होते हैं।

एक अत्यंत सरल उदाहरण के रूप में पूर्वक्रमों पर विचार करें। पूर्वक्रम ऐसा
संबंध~$R$ है जो किसी समुच्चय~$A$ पर स्वतुल्य और संक्रामक दोनों है। समुच्चय
~$A$ पर द्वि-स्थानीय संबंध $R \subseteq A \times A$ ठीक वही सामग्री है जो
एक ही द्वि-स्थानीय संबंध-चिह्न~$\Obj P$ वाली प्रथम-क्रम भाषा के लिए
!!a{structure} देने को चाहिए: हम $\Domain{M} = A$ और
$\Assign{\Obj P}{M} = R$ रखेंगे। चूँकि $R$~एक पूर्वक्रम है, वह स्वतुल्य और
संक्रामक है, और हम प्रथम-क्रम तर्क के !!{sentence}s का ऐसा समुच्चय~$\Gamma$
पा सकते हैं जो यही कहता है:
\begin{align*}
  & \lforall[\Obj v_0][\Atom{\Obj P}{\Obj v_0,\Obj v_0}]\\
  & \lforall[\Obj v_0][\lforall[\Obj v_1][\lforall[\Obj v_2][((\Atom{\Obj P}{\Obj v_0,\Obj v_1} \land \Atom{\Obj P}{\Obj v_1,\Obj v_2}) \lif \Atom{\Obj P}{\Obj v_0,\Obj v_2})]]]
\end{align*}
!!{sentence}s के ये दो उदाहरण केवल ``किसी भी~$x$ के लिए $Rxx$'' ($R$ स्वतुल्य है) और
``जब भी $Rxy$ तथा $Ryz$, तब $Rxz$ भी'' ($R$ संक्रामक है) के प्रतीकीकरण हैं।
हम देखते हैं कि !!a{structure}~$\Struct{M}$ इन दो !!{sentence}s से बने~$\Gamma$ का
प्रतिरूप तभी और केवल तभी है जब $R$ (अर्थात्~$\Assign{\Obj P}{M}$),~$A$ पर
पूर्वक्रम हो (अर्थात्~$\Domain{M}$ पर)। दूसरे शब्दों में,~$\Gamma$ के
प्रतिरूप ठीक सभी पूर्वक्रम हैं। सभी पूर्वक्रमों का कोई भी ऐसा गुण जिसे केवल
~$\Obj P$ को !!{predicate} मानने वाली प्रथम-क्रम भाषा में व्यक्त किया जा सके
(जैसे ऊपर स्वतुल्यता और संक्रामकता),~$\Gamma$ के दो !!{sentence}s से निष्पन्न
होता है, और इसका विलोम भी सत्य है। अतः~$\Gamma$ के प्रतिरूपों के बारे में हम
जो कुछ सिद्ध कर सकते हैं, वह सभी पूर्वक्रमों के बारे में सिद्ध कर चुके हैं।

किसी विशिष्ट सिद्धान्त और प्रतिरूप-वर्ग (जैसे~$\Gamma$ और सभी पूर्वक्रम) के
लिए इस बारे में रोचक प्रश्न होंगे कि संगत प्रथम-क्रम भाषा में क्या व्यक्त
किया जा सकता है और क्या नहीं। !!{structure}s कुछ ऐसे गुण रखती हैं जो सभी
भाषाओं और प्रतिरूप-वर्गों के लिए रोचक हैं, विशेषतः !!{domain} के आकार से
संबंधित गुण। उदाहरणार्थ, यह सदा व्यक्त किया जा सकता है कि !!{domain} में ठीक
$n$~!!{element}s हैं, किसी भी~$n \in \PosInt$ के लिए। अनंत रूप से अनेक
!!{sentence}s के किसी समुच्चय से यह भी व्यक्त किया जा सकता है कि !!{domain}
अनंत है। किंतु यह व्यक्त नहीं किया जा सकता कि प्रांत परिमित है, अथवा यह कि
प्रांत !!{nonenumerable} है। प्रथम-क्रम भाषाओं की सीमाओं संबंधी ये तथ्य
संहतता और लोवेनहाइम--स्कोलेम (L\"owenheim--Skolem) प्रमेयों से निकलते हैं।

\end{document}
